\documentclass[10pt,a4paper]{amsart}
\usepackage[margin=19mm]{geometry}
\usepackage{amsmath,amssymb,mathtools}
\usepackage[scr=rsfs,cal=euler]{mathalfa}
\usepackage{xfrac}
\usepackage[unicode,hidelinks,bookmarks=false]{hyperref}
\newcommand{\ie}{i.e.\ }
\newcommand{\cf}{cf.\ }
\newcommand{\resp}{respectively\ }
\newcommand{\Subsection}{\subsection}
\providecommand{\nobreakdash}{\nobreak\hbox{-}}
\setlength{\emergencystretch}{2em}
\multlinegap0pt
\usepackage{bm}
\usepackage{bbm}
\usepackage{dsfont}
\usepackage{xspace}
\usepackage{cancel}
\usepackage{mathtools}
\usepackage{amsthm}
\makeatletter
\@ifundefined{theorem}{\newtheorem{theorem}{Theorem}}{}
\makeatother
\title[On the true detection probability of the uniformly optimal s]{On the true detection probability of the uniformly optimal search plan}
\author{Liang Hong}
\date{Source: arXiv:2311.18226v4 (CC BY 4.0)}
\begin{document}
\maketitle

\noindent\emph{Source and licence.} On the true detection probability of the uniformly optimal search plan. Version arXiv:2311.18226v4 (CC BY 4.0), distributed under the Creative Commons Attribution 4.0 International licence, as recorded in the arXiv licence metadata for this exact version and shown on the abstract page https://arxiv.org/abs/2311.18226v4. Full rights notice: licenses/arxiv-2311.18226-CC-BY-4.0.txt.\par\smallskip

\noindent\textit{Editorial scope.} Counted unit: the Theorem whose source label is \texttt{thm:nooptdiscrete1} in the article (source0.tex lines 486--494, source label \texttt{thm:nooptdiscrete1}), delivered together with the article's own complete original proof of it (source0.tex lines 497--593, one \texttt{proof} environment of 97 lines). No mathematics is written, repaired or re-derived: statement and proof are the article's own text line for line. Required context retained uncounted: eq:trueprob (lines 173--176); eq:truemean (lines 200--203); eq:rate (lines 241--244). Every definition, display and label of those lines is retained unchanged. Omitted from this package and not redistributed: 1--172, 177--199, 204--240, 245--485, 495--496, 594--857. Nothing inside a retained excerpt is omitted or altered: every retained body is delivered exactly as the article's lines are written, so no omission receipt of any kind accompanies this package. Citations: the retained excerpts carry no citation command at all, so no bibliography entry of the article is transcribed and no reference is redistributed. Reference adaptation: none was needed and none was made; every reference inside the delivered unit resolves inside it, so no unresolved reference or citation is produced. Printed slips kept literal: no printed word, symbol, hypothesis or number of the counted statement or of its proof was altered. Layout and class: the article's own class is \texttt{article} with packages including color, graphicx, subfigure, amsmath, amsthm, amssymb, amsfonts, fullpage, ifthen, url, natbib, multirow, bm, tikz-qtree; the wrapper is the lane's pinned \texttt{amsart} editorial wrapper at 10pt on A4, which restates the theorem-like environments the retained text needs and adds the article's own macro definitions that the retained text needs (none), with extra wrapper packages (bm, bbm, dsfont, xspace, cancel, mathtools). The delivered numbering is the unit's own. Assets: no retained span contains a figure environment, a graphics-inclusion command, a PDF-page inclusion, a table or a TikZ picture; no image file is copied into this package, the package declares no asset dependency and no image file is redistributed with it. Licence: version arXiv:2311.18226v4 of this article is distributed under the Creative Commons Attribution 4.0 International licence, as recorded in the arXiv licence metadata for this exact version and shown on the abstract page https://arxiv.org/abs/2311.18226v4. A full rights notice is delivered at licenses/arxiv-2311.18226-CC-BY-4.0.txt. Characters: every retained excerpt is pure ASCII, so the delivered engine, XeLaTeX with its default Latin Modern text font, reports no missing character.
\begin{equation}
\label{eq:trueprob}
P^{\#}[f]=d(x_0, f(x_0)).
\end{equation}

\begin{equation}
\label{eq:truemean}
\mu^{\#} (\varphi)=\int_0^\infty(1-P^{\#}[\varphi(\cdot, t)])dt.
\end{equation}

\begin{equation}
\label{eq:rate}
q_x (y)=\pi(x)\frac{\partial }{\partial y}d(x, y), \quad x\in X \text{ and $y\geq 0$},
\end{equation}

\begin{theorem}
\label{thm:nooptdiscrete1}
Assume $X$ is discrete and $x_0$ is the true cell in which the target is located.  Suppose the cost function takes the form $c(x,y)=y$ for all $y\geq 0$ and $x \in X$,  the detection function is regular,  $\Pi$ is a target distribution on $X$, and $\varphi^\star$ is the uniformly optimal search plan for $\Pi$ within $\Phi(E)$ such that $0\leq \varphi^\star(x_0, t)<E(t)$.  If either of the following conditions holds:
\begin{enumerate}
\item[(i)]there is an $x_1\in X$ such that $x_1\neq x_0$ and $\pi(x_1)>\pi(x_0)$, 
\item[(ii)]$1>\pi(x_0)\geq \pi(x)$ for all $x\neq  x_0$ and  $\ E(t)>q_{x_1}^{-1}(q_{x_2}(0))$ for all $t>0$,
\end{enumerate}
then there exists another target distribution $\widetilde{\Pi}\neq \Pi$ such that $P^{\#}[\widetilde{\varphi}^\star(\cdot, t)]>P^{\#}[\varphi^\star(\cdot, t)]$ for all $t>0$ and $\mu^{\#}(\widetilde{\varphi}^\star)<\mu^{\#}(\varphi^\star)$, where $\widetilde{\varphi}^\star$ is the uniformly optimal search plan for $\widetilde{\Pi}$ within $\Phi(E)$.
\end{theorem}

\begin{proof}
The inverse function of the  rate of return $q_x(y)=\pi(x) \frac{\partial}{\partial y}d(x,y)$ is given by
\begin{equation*}
q^{-1}_x(\lambda)=\left\{
		                           \begin{array}{ll}
		                           \text{the inverse function of $q_x(y)$  evaluated at $\lambda$},& \hbox{if $0<\lambda\leq q_x(0)$,} \\
					0, & \hbox{if $\lambda>q_x(0)$.} 
		                          \end{array}
		                         \right.
\end{equation*}	
Suppose $\widetilde{\pi}$ is the new target density to be constructed.  Since $\frac{\partial}{\partial y}d(x, y)$ is independent of the target density $\pi$,  $\widetilde{\pi}(x_0)>\pi(x_0)$ will imply $\widetilde{q}_{x_0}(y)> q_{x_0}(y)$ for all $y\geq 0$.  On the other hand,  $q_x^{-1}(\lambda)$ equals the inverse of $\ \frac{\partial}{\partial y}d(x, y)$ evaluated at $\frac{\lambda}{\pi(x)}$ when $\lambda\leq q_x(0)$.  Since $\ \frac{\partial}{\partial y}d(x, y)$ is strictly decreasing, we have
\begin{equation}
\label{eq:discrete1}
\widetilde{q}^{-1}_{x_0}(\lambda)\left\{
		                           \begin{array}{ll}
		                           >q_{x_0}^{-1}(\lambda),& \hbox{for $0<\lambda\leq \widetilde{q}_{x_0}(0)$,} \\
					=q_{x_0}^{-1}(\lambda)=0, & \hbox{if $\lambda>\widetilde{q}_{x_0}(0)$.} 
		                          \end{array}
		                         \right.
\end{equation}	
Without loss of generality, we may assume $\widetilde{Q}^{-1}(E(t))\leq \widetilde{q}_{x_0}(0)$.
We are going to construct a new target distribution $\widetilde{\Pi}$ such that
\begin{equation}
\label{eq:condition}
\widetilde{q}_{x_0}^{-1}(\widetilde{Q}^{-1}(E(t)))>q_{x_0}^{-1}(Q^{-1}(E(t))),
\end{equation}
which will imply $P^{\#}[\widetilde{\varphi}^\star(\cdot, t)]>P^{\#}[\varphi^\star(\cdot, t)]$ for all $t>0$ and $\mu^{\#}(\widetilde{\varphi}^\star)<\mu^{\#}(\varphi^\star)$ via (\ref{eq:trueprob}) and (\ref{eq:truemean}). To this end,  consider two cases: (i)~there is an $x_1\in X$ such that $x_1\neq x_0$ and $\pi(x_1)>\pi(x_0)$ and (ii)~$1>\pi(x_0)\geq \pi(x)$ for all $x\neq  x_0$ and $E(t)>q_{x_1}^{-1}(q_{x_2}(0))$ for all $t>0$.

In Case~(i), we construct the new target density $\widetilde{\pi}$ by switching the probability masses of $\pi$ on $x_0$ and $x_1$ and keep everything else intact. That is, we define
\[ \widetilde{\pi}(x)=\left\{
		                           \begin{array}{ll}
		                           \pi(x_1),& \hbox{if $x=x_0$,} \\
					\pi(x_0), & \hbox{if $x=x_1$,} \\
					\pi(x), & \hbox{otherwise.}
		                          \end{array}
		                         \right.		                         
		                         \]
Then $\widetilde{\pi}(x_0)>\widetilde{\pi}(x_1)$.  By (\ref{eq:discrete1}), we have $\widetilde{q}^{-1}_{x_0}(\lambda)>q^{-1}_{x_0}(\lambda)$ for $0<\lambda\leq \widetilde{q}_{x_0}(0)$.  Also,   the definition of $\widetilde{\pi}$ and (\ref{eq:rate}) imply 
\[\widetilde{Q}(\lambda)=\sum_{x\in X}\widetilde{q}_x^{-1}(\lambda)=\sum_{x\in X}q_x^{-1}(\lambda)=Q(\lambda), \quad \text{for all $\lambda>0$},\]
which further implies $\widetilde{Q}^{-1}(K)=Q^{-1}(K)$ for all  $K>0$. Therefore, 
\[\widetilde{q}^{-1}_{x_0}(\widetilde{Q}^{-1}(E(t)))=\widetilde{q}^{-1}_{x_0}(Q^{-1}(E(t)))>q_{x_0}^{-1}(Q^{-1}(E(t))).\]

In Case~(ii), we may assume without loss of generality that $X=\{1, \ldots, m\}$ for some integer $m\geq 2$, $x_0=1$, and $1>\pi(1)>\pi(2)\geq \pi(3)\geq \ldots \geq \pi(m)>0$.  Then the definition of $q_x$ implies $q_1(0)>q_2(0)\geq q_3(0)\ldots\geq q_m(0)$. For $i=1, \ldots, m$, we have
\begin{equation*}
q^{-1}_i(\lambda)=\left\{
		                           \begin{array}{ll}
		                           \left(\frac{\partial}{\partial y}d(i, y)\right)^{-1}\bigg |_{y=\frac{\lambda}{\pi(i)}},& \hbox{if $0<\lambda\leq q_x(0)$,} \\
					0, & \hbox{if $\lambda>q_i(0)$.} 
		                          \end{array}
		                         \right.
\end{equation*}	
Thus, 
\[
Q(\lambda)=\sum_{i=1}^m q_i^{-1}(\lambda)=\left\{
		                           \begin{array}{ll}
		                           \sum_{i=1}^m \left(\frac{\partial}{\partial y}d(i, y)\right)^{-1}\bigg |_{y=\frac{\lambda}{\pi(i)}}, & \hbox{if $0<\lambda\leq q_m(0)$,} \\
					\sum_{i=1}^{m-1} \left(\frac{\partial}{\partial y}d(i, y)\right)^{-1}\bigg |_{y=\frac{\lambda}{\pi(i)}}, & \hbox{if $q_m(0)<\lambda\leq q_{m-1}(0)$,} \\
						\ldots & \hbox{$\ldots$,} \\
 \left(\frac{\partial}{\partial y}d(1, y)\right)^{-1}\bigg |_{y=\frac{\lambda}{\pi(1)}} + \left(\frac{\partial}{\partial y}d(2, y)\right)^{-1}\bigg |_{y=\frac{\lambda}{\pi(2)}}, & \hbox{if $q_3(0)<\lambda\leq q_2(0)$,} \\
 \left(\frac{\partial}{\partial y}d(1, y)\right)^{-1}\bigg |_{y=\frac{\lambda}{\pi(1)}}, & \hbox{if $q_2(0)<\lambda\leq q_1(0)$,} \\
                 0, & \hbox{if $\lambda>q_1(0)$.} 
		                          \end{array}
		                         \right.
\]
It follows that $Q^{-1}$ has the form
\[
Q^{-1}(K)=\left\{
		                           \begin{array}{ll}
		                          q_1(K),& \hbox{if $0<K \leq q_1^{-1}(q_2(0))$,} \\
		                         (q_1^{-1}+q_2^{-1})^{-1}(K),& \hbox{if $q_1^{-1}(q_2(0))<K\leq (q_1^{-1}+q_2^{-1})(q_3(0))$,} \\
		                          \ldots,& \hbox{if $\ldots$},
		                          \end{array}
		                         \right.
\]
and $\varphi^\star$ is of the form
\begin{eqnarray*}
\varphi^\star(1, t) &=& q_1^{-1}(Q^{-1}(E(t)))\\
&=&\left\{
		                           \begin{array}{ll}
		                         E(t),& \hbox{if $0<E(t) \leq q_1^{-1}(q_2(0))$,} \\		    
		                         q_1^{-1}((q_1^{-1}+q_2^{-1})^{-1}(E(t))), & \hbox{if $q_1^{-1}(q_2(0))< E(t) \leq (q_1^{-1}+q_2^{-1})(q_3(0))$,} \\	                 
		                          \ldots,& \hbox{if $\ldots$},
		                          \end{array}
		                         \right.
\end{eqnarray*}
Consider the case $q_1^{-1}(q_2(0))< E(t) \leq (q_1^{-1}+q_2^{-1})(q_3(0))$. By the definition of $q_x^{-1}$, we know 
$\varphi^\star$ is continuous in both $\pi(1)$ and $\pi(2)$ and 
\[
\lim_{\pi(1)\rightarrow 1^-} \varphi^\star(1, t)=\lim_{\pi(1)\rightarrow 1^-}q_1^{-1}((q_1^{-1}+q_2^{-1})^{-1}(E(t)))=q_1^{-1}((q_1^{-1})^{-1}(E(t)))=E(t).
\]
For all other cases where $E(t)>(q_1^{-1}+q_2^{-1})(q_3(0))$,  a similar argument leads to the same conclusion. 
Therefore, there is another target density function $\widetilde{\pi}$ such that (\ref{eq:condition}) holds as far as $E(t)>q_1^{-1}(q_2(0))$ for all $t>0$.




\end{proof}

\end{document}
